% SUPERSEDED: registered audited version is sections/A_compact_type_fixed.tex.
% UNREGISTERED A.41 candidate, prepared for parent audit during checkpoint freeze.
% The proposed Proved status must not be added to the register before that audit.
\subsection{A compact-type criterion for higher-genus covers}

\begin{statement}{A.41}{Proved}{A rational marked path in compact-type reduction forces cusp fixed groups}
Let \(X/k\) be a smooth affine hyperbolic curve over a number
field, whose smooth compactification has positive genus and
whose complement contains at least two geometric points.
Let \(P\) be a rational cusp and let \(s\) be any packet
at \(P\). Suppose that, after a finite extension at a place
of residue characteristic \(p\), the pointed stable model
has a tree as its dual graph. Suppose also that a distinct
cusp \(Q\) has this property: the unique path between the
vertices carrying \(P\) and \(Q\), including its endpoints,
contains only genus-zero components. Then, for every prime
\(\ell\ne p\) and every open subgroup
\(\Gamma_0\subset\overline{[G_k,G_k]}\),
\[
 \left(\pi_1(X_{\bar k})^{(\ell)}\right)^{s(\Gamma_0)}
                         =I_{P,\ell}.                 \tag{A.41.1}
\]
The same conclusion holds for the permitted open
\(\Delta_0\) in Part A whenever the cover and selected
cusp are preserved and a finite field of definition has
its closed derived subgroup contained in \(\Delta_0\).

A sufficient condition for the geometric hypothesis is
compact-type reduction with a special Jacobian of
\(p\)-rank zero and a \(p\)-primary torsion cusp divisor
\([P-Q]\). More generally it suffices that \(P,Q\) have
equal specialization under the extended Abel--Jacobi map.
This gives a criterion for higher-genus instances of
Task A.5; it makes no assertion for a dual graph with a
cycle, or for a path which contains a positive-genus vertex.
\end{statement}

\begin{proof}
\paragraph{Field restriction and the local graph.}
Enlarge the number field to \(L\) so that all cusps and
components are split, the stated pointed stable model is
defined at a place above \(p\), and the Zariski closure of
the proper \(\ell\)-adic homology representation is
connected. We may also arrange
\(\Gamma=\overline{[G_L,G_L]}\subset\Gamma_0\), or
\(\Gamma\subset\Delta_0\) in Part A. Indeed an open
subgroup of \(G_k\) can be chosen whose intersection
with the closed derived group lies in the prescribed open
subgroup, and its own closed derived group lies in that
intersection. Further finite extensions preserve the graph
hypothesis: edge lengths are rescaled and any inserted
valence-two rational vertices may be contracted. It is
enough to prove that the actual \(\Gamma\)-fixed group
is \(I_{P,\ell}\). The specified larger group already
fixes this inertia, since its cyclotomic character is one.

Let \(C_1,\ldots,C_m\) be the positive-genus components
of the compact-type special fibre. The good reduction of
its Jacobian and its special product decomposition give
an orthogonal symplectic decomposition
\[
 V=H_1(\overline X_{\bar L},\mathbf Q_\ell)
          =\bigoplus_{i=1}^m V_i,
 \qquad \Omega=\sum_{i=1}^m\Omega_i,
 \qquad \Omega_i\in\Lambda^2V_i.                    \tag{A.41.2}
\]
Here \(\Omega_i\) is the nondegenerate symplectic
bivector on the indicated component. This is a local
symplectic decomposition; the individual summands are not
asserted to be globally Galois stable. The compact-type
Jacobian product is recalled in
\cite[Section 1, p.1]{fghDupuyRabinoff}. Equivalently,
(A.41.2) is the homology decomposition obtained by cutting
the surface along the separating nodes in the based
comparison of \cite[Theorem 3.1.6]{fghBettsDogra}.

Root the metric tree at the vertex carrying \(P\), and
orient its finite edges away from that root. Write
\(m_e>0\) for their lengths. For each edge let
\(i_e,Q_e,j_e\) indicate whether component \(C_i\),
the cusp \(Q\), and the cusp \(j\), respectively,
lie on its side away from the root. Use handle generators
of weight one for \(V\), and cusp generators \(t_j\)
of weight two. The first peripheral relation is
\(\Omega+\sum_j t_j=0\). Eliminate \(t_Q\) and
retain \(t_P\) among the independent cusp generators.
The weighted graded geometric Lie algebra is then
\[
                         \mathbb L(V\oplus T),
 \qquad T=\langle t_j:j\ne Q\rangle,
 \quad \operatorname{wt}V=1,\quad\operatorname{wt}T=2.
                                                               \tag{A.41.3}
\]

\paragraph{The genuine weighted derivation.}
Before eliminating \(Q\), put
\[
 \delta_e=\sum_i i_e\Omega_i+\sum_{j:\,j_e=1}t_j,
 \qquad W_i=\sum_e m_ei_e\delta_e,
 \qquad W_j=\sum_e m_ej_e\delta_e.                    \tag{A.41.4}
\]
The actual based edge-twist formula of
\cite[Definition 3.1.5 and Theorem 3.1.6]{fghBettsDogra}
gives the first weighted monodromy derivation as
\[
                  Dv=[W_i,v]\ (v\in V_i),
                  \qquad Dt_j=[W_j,t_j].             \tag{A.41.5}
\]
As in A.36 and A.40, choose the ordinary comparison point
in the actual selected cusp disc. Comparison with the
specified tangential packet changes the first derivation
only by \(\lambda\operatorname{ad}(t_P)\), with
\(\lambda\in\mathbf Q_\ell\). Thus in (A.41.4)--(A.41.5)
replace every \(W_i,W_j\) by \(W_i+\lambda t_P,W_j+
\lambda t_P\). There is no weight-one inner discrepancy:
it would centralize the independent cusp generator
\(t_P\), whose Lie centralizer has no weight-one part.
A common change of edge orientation merely negates \(D\).

After semistable extension, the prime-to-\(p\) monodromy
has this actual edge action. A local Frobenius--tame-inertia
commutator, regarded inside \(\Gamma\), has first
weighted derivation a nonzero common multiple of \(D\),
with factor \(q_{\mathrm{res}}-1\ne0\). Thus the
leading classes of the unipotent closure of actual
\(\Gamma\)-fixed elements are killed by \(D\).
All terms in (A.41.5) shift weight by two, and the cusp
packet term has exactly the same shift.

The coefficient of \(\Omega_i\) in the eliminated
\(\delta_e\) is \(i_e-Q_e\). Its coefficient in
its own \(W_i\) is therefore
\[
               r_i=\sum_e m_ei_e(1-Q_e)>0.           \tag{A.41.6}
\]
The equality describes the distance from the branching
point of \(P,Q,C_i\) to \(C_i\). Since the entire
\(P\)--\(Q\) path has genus zero, \(C_i\) is off
that path and the indicated distance is strictly positive.
This is the only positivity assertion needed. Coefficients
of \(\Omega_i\) in other \(W\)'s may be negative,
zero, or equal to \(r_i\), and \(\lambda\) is arbitrary.

\paragraph{An unbounded calculation with several handle blocks.}
Work over \(F=\mathbf Q_\ell\) in the free associative
enveloping algebra \(A=T(V\oplus T)\), and put
\[
             U=\langle\Omega_1,\ldots,\Omega_m\rangle,
             \qquad B=F\langle T,U\rangle\subset A.
                                                               \tag{A.41.7}
\]
This is a free associative algebra on the displayed
primitive generators. To make its complementary basis
explicit, choose a symplectic basis in each \(V_i\),
with leading word \(a_{i1}b_{i1}\) for \(\Omega_i\).
Its coefficient is one and it has no proper self-overlap.
Each handle algebra has the triangular basis
\[
 \{\Omega_i^n:n\ge0\}\ \sqcup
 \{\Omega_i^u w\Omega_i^v:u,v\ge0,
   \ w\ne\varnothing\text{ has neither initial nor terminal }
                                      a_{i1}b_{i1}\}.           \tag{A.41.8}
\]
Indeed strip the maximal initial and terminal powers of
\(a_{i1}b_{i1}\) from every word; the resulting distinct
leading words exhaust the usual word basis. Taking the
free product of these component algebras and the free
cusp algebra gives a normal basis in which every word
outside \(B\) has a first nonboundary handle core,
of some type \(i\), preceded by an arbitrary \(B\)-word.
Pure boundary powers adjoining that core are retained
in the adjoining slots. A separator from a different
component or a cusp letter cannot disappear.

By (A.41.5), the derivation on every handle word of
component type \(i\) is exactly its commutator with
\(W_i\). Also
\(D\Omega_i=[W_i,\Omega_i]\), so \(D\) preserves
\(B\); write this restriction as \(D_B\). For a
fixed first core of type \(i\), the part increasing
the length of its first boundary slot is
\[
                             K_i=D_B+R_{W_i}.        \tag{A.41.9}
\]
All terms acting after that core preserve the first-slot
length. They cannot remove its core: multiplication by
\(\Omega_i\) at either end stays in its displayed
bimodule, whereas a different boundary letter adds a
separator. The derivative of an existing boundary letter
retains that letter in both commutator terms. Thus the
top first-slot operator acts independently for each fixed
first-core tail, with no contribution from later slots.

To prove \(K_i\) injective, count occurrences of the
single letter \(\Omega_i\) in \(B\). Its part which
raises that count by one is
\[
                K_i^+=D_{c_i}+r_iR_{\Omega_i},
 \quad D_{c_i}\Omega_i=0,
 \quad D_{c_i}u=c_u[\Omega_i,u]\quad(u\ne\Omega_i), 
                                                               \tag{A.41.10}
\]
where \(c_u\) is the coefficient of \(\Omega_i\)
in the \(W\) belonging to the boundary letter \(u\).
For \(u=\Omega_j\) use \(W_j\) for component \(j\);
for a cusp use its cusp \(W\). The self coefficient
in \(D_B\Omega_i=[W_i,\Omega_i]\) vanishes, so
that expression preserves this count. Every omitted
term in \(K_i\) also preserves it.

Fix a word \(u_1\cdots u_n\) in all the boundary
letters other than \(\Omega_i\). Insert powers of
\(\Omega_i\) in its \(n+1\) slots and identify the
resulting vector space with \(F[Z_0,\ldots,Z_n]\).
The operator (A.41.10) becomes multiplication by
\[
 c_{u_1}Z_0+\sum_{h=1}^{n-1}(c_{u_{h+1}}-c_{u_h})Z_h
                           +(r_i-c_{u_n})Z_n.        \tag{A.41.11}
\]
Its coefficients telescope to \(r_i\ne0\). When
\(n=0\), its multiplier is \(r_iZ_0\). Therefore
it is injective on every mutually independent word
summand. Taking the largest \(\Omega_i\)-count of
a finite nonzero input proves that \(K_i\) is
injective. Now take a largest first-slot length among
the normal words of a nonzero polynomial outside
\(B\). Formula (A.41.9) cannot vanish there. Hence,
in every weight,
\[
                   \ker D\subset B\cap\mathbb L(V\oplus T)
                                      =\mathbb L(T\oplus U).
                                                               \tag{A.41.12}
\]
The final equality holds because \(B\) is the tensor
Hopf algebra on the indicated freely generated primitive
elements. This argument proves the assertion in all
weights, and does not infer it from finitely many ranks.

\paragraph{The exact intersection of the global translates.}
Let \(\mathfrak h\) be the Lie algebra of the
unipotent closure of the actual \(\Gamma\)-fixed
subgroup, with its induced weight filtration, as in
A.38. Its leading graded pieces are stable under
\(G_L\) and, by (A.41.12), lie in
\(\mathbb L(T\oplus U)\). On the weighted ambient
grading, \(V\) is the canonical proper homology,
\(T\) is the independent boundary subspace with
character \(\chi_\ell\), and \(\Lambda^2V\) is
the canonical space of brackets of weight-one classes.
In particular global Galois action does not mix \(T\)
with \(\Lambda^2V\) in weight two.

Let \(G\subset\operatorname{GSp}(V)\) be the
connected Zariski closure of that action. Each finite
weighted piece of \(\mathbb L(V\oplus T)\) is an
algebraic representation of \(G\), with the action on
\(T\) given by its similitude character. The same
leading piece of \(\mathfrak h\) therefore lies in
all translates \(\mathbb L(T\oplus gU)\),
\(g\in G\). Set
\[
                           U_* =\bigcap_{g\in G}gU.
                                                               \tag{A.41.13}
\]
The following intersection is an equality of actual
subspaces of the free weighted Lie algebra:
\[
        \bigcap_{g\in G}\mathbb L(T\oplus gU)
                                    =\mathbb L(T\oplus U_*).
                                                               \tag{A.41.14}
\]
Here is a direct verification. The block map
\[
          T\bigl(T\oplus(V\otimes V)\bigr)
                               \longrightarrow T(V\oplus T)
\]
is injective: cusp letters are individually recognizable,
and each maximal string of handle letters in its image
has even length, with a unique consecutive pairing from
its left end. Thus, in weight \(2n\), intersections
can be taken inside the \(n\)-fold tensor powers of
\(T\oplus(V\otimes V)\). For subspaces \(A_1,A_2\)
of a vector space, a basis adapted to their intersection
gives
\(A_1^{\otimes n}\cap A_2^{\otimes n}
=(A_1\cap A_2)^{\otimes n}\). Iterate this identity.
The intersection of any family of finite-dimensional
subspaces is already a finite subintersection, so the
same assertion holds for the family in (A.41.13).
Taking primitive elements gives (A.41.14). No passage
through a surface-group quotient is involved.

\paragraph{Why the remaining component classes are Tate.}
Use the total nondegenerate bivector \(\Omega\) to
identify a bivector \(\beta:V^*\to V\) with
\(A_\beta=\beta\Omega^{-1}\in\operatorname{End}(V)\).
Then \(A_{\Omega_i}\) is the projector onto the
orthogonal block \(V_i\). Consequently \(U\)
identifies with the commutative diagonal algebra
\[
                 \mathcal D=\bigoplus_{i=1}^mF P_i,
          \qquad P_i|_{V_i}=1,\quad P_i|_{V_j}=0\ (j\ne i).
                                                               \tag{A.41.15}
\]
For a symplectic similitude \(g\),
\[
               A_{g\beta}=\chi_\ell(g)\,gA_\beta g^{-1}.
                                                               \tag{A.41.16}
\]
The subspace corresponding to \(U_*\) is conjugation
stable and lies in \(\mathcal D\). For any element
\(A\) in that subspace, its entire \(G\)-conjugacy
orbit is contained in \(\mathcal D\) and has the
same multiset of eigenvalues as \(A\). There are
only finitely many elements of \(\mathcal D\) with
that multiset: each fixed nonzero block must receive
one of the finitely many eigenvalues. This remains
true with repeated eigenvalues or unequal block sizes.
After extension to an algebraic closure the orbit of
the connected group \(G\) is connected, and so this
finite orbit is a singleton. Hence \(G\) centralizes
\(A\). Formula (A.41.16) proves
\[
                         g\beta=\chi_\ell(g)\beta
                      \quad\text{for all }\beta\in U_*.
                                                               \tag{A.41.17}
\]
No reductivity assumption is used, and none of the
individual local projectors \(P_i\) is claimed to
come from a global endomorphism.

Every generator in \(T\oplus U_*\) now has weight
two and character \(\chi_\ell\). Equations
(A.41.12)--(A.41.17) show that \(\mathfrak h\) has
zero odd graded pieces and scalar character
\(\chi_\ell^n\) in weight \(2n\). Apply the
actual-subgroup criterion in A.38. It gives that the
\(\Gamma\)-fixed group is exactly \(I_{P,\ell}\),
which proves (A.41.1) and its stated Part A variant.

\paragraph{A torsion-divisor criterion for the rational path.}
The extended Abel--Jacobi map exists on the pointed
compact-type model by
\cite[Theorem 1.1]{fghDupuyRabinoff}. On any component,
its projection to that component's Jacobian is its
Abel map, up to translation, and all other component
coordinates are constant; this is
\cite[Proposition 3.2, pp.5--6]{fghDupuyRabinoff}.
It follows that two distinct smooth marked points with
the same Abel specialization cannot have a positive-genus
component on their joining path. Indeed, projection to
the Jacobian of such a component gives the difference
of the two points at which the path enters and leaves
it, using the marked endpoint if necessary. Those
points are distinct: two different incident edges have
different nodes, and a marked point is disjoint from
all nodes. The Abel map of a smooth positive-genus
component is injective on points, so this projected
difference is nonzero, a contradiction. Thus their
path contains only rational components.

Finally suppose the special Jacobian has \(p\)-rank
zero and \([P-Q]\) is \(p\)-primary torsion. Its
specialization is a geometric \(p\)-power torsion
point of that abelian variety, hence is zero. The
Abel specializations are therefore equal, giving the
sufficient condition in the statement.

The exact program \texttt{code/a\_compact\_type\_fixed.py}
checks the rooted-tree distance identity, multicore
prefix operators with negative and repeated coefficients,
and tensor intersections for a symplectic mixing matrix.
The all-degree and connected-group arguments are proved
above, independently of those finite certificates.
\end{proof}
